\section{Discussion}

\subsection{Analysis of Findings}

We analyze failure cases in integration tests and benchmarking for deeper insight into reasons for inaccurate answers.

\paragraph{Document Comprehension}
We find that document layout comprehension plays an important role in retrieval accuracy.
Documents tested in benchmarks include real-world documents from a wide range of domains with various formats.
For instance, presentation slides, academic papers with multi-column layouts, and informational leaflets or brochures
that feature interleaved text in images occupying the entire page space.
Our document processor extracts images and texts in reading order, hence possibly confusing multi-column document layouts for single-column layouts.
While our pipeline is observed to successfully retrieve entire presentation slides to produce a correct answer,
we also find that in some cases, retrieved raw text documents show content in the same line among multiple columns.
Such noise may prevent correct generation of answers, causing failure cases.

\paragraph{Knowledge Graph Storage}
Our graph-based storage option extracts entities and relationships to form a knowledge graph.
We then generate community summaries with an LLM for communities of entity nodes detected in the graph.
Theoretically graph-based storage preserves long-form context more accurately;
however in our tests we find that flattening source documents into summaries causes loss of specific details in source documents.
This led to a gap between dense retrieval and graph community retrieval for augmenting generation.
We hence default to using dense retrieval.

\paragraph{Generation Accuracy}
The level of generation accuracy increases from zero-shot to text-only retrieval, then to multimodal retrieval.
Inline images may contain valuable fine-grained visual information useful in providing an answer to a query.
By retrieving related images and passing them to the MLLM, we make full use of external timely knowledge to support generation grounding,
increasing the accuracy of final answers.
This is observed in our end-to-end tests, where the highest accuracy score obtained from multimodal retrieval is 54.6\%,
and that of text-only retrieval only reaches 45.5\%.
Zero-shot generation gives 27.3\% accuracy.
Our findings verify our hypothesis that retrieving visual information helps in reducing hallucination and factual errors in answering queries.


\subsection{Limitations}

Despite our proposed pipeline's strengths, we identify three major limitations: 
preprocessing efficiency, theoretical bounds of dense retrieval, and current MLLM visual perception capability issues.

\paragraph{Preprocessing}
The preprocessing step in our pipeline requires inline image extraction from documents and captioning.
We adopt MLLM image captioning with source document context as this method gives the highest accuracy among all methods tested.
While the accuracy is exceptional (reaching an average of 97.79\% on our small test corpus),
we find that processing time reaches 3.89 seconds per image.
This creates a large indexing overhead for real-world uses, especially corporate-grade document batches that may contain thousands or millions of images.
Using proprietary closed-source MLLMs via API requests might incur further networking or connection delays,
while running inference on open-source MLLMs in local environments would demand much computation power and hence more powerful hardware.

Our proposed pipeline also does not support video processing.
Similar work such as VideoRAG \cite{ren2025videoragretrievalaugmentedgenerationextreme} already demonstrate the possibility of 
indexing long-form videos into a knowledge base (a knowledge graph in this specific implementation) and can be used for generation grounding.
As videos are formed of continuous frames of images, it would be possible to adopt a similar approach to our pipeline in processing videos.

\paragraph{Dense Retrieval Bounds}
Our pipeline offers dense retrieval from a vector database or knowledge graph retrieval for the text-to-text retrieval step,
yet dense retrieval in the DINOv2 vector space is the only method for text-to-image and image-to-image retrieval using the DINOv2 representation space.
Although dense retrieval is extensively used in RAG applications \cite{lewis2021retrievalaugmentedgenerationknowledgeintensivenlp,tonmoy2024comprehensivesurveyhallucinationmitigation},
recent research has revealed the theoretical limits of dense retrieval.
\cite{weller2026theoreticallimitationsembeddingbasedretrieval} finds that dense retrieval faces theoretical limits of retrieval effectiveness
with regards to the total number of documents $n$, the number of documents to retrieve $k$, and the dimensionality of dense vectors used in indexing documents.
In corporate-level applications of our pipeline, where the number of documents and images can reach large numbers, theoretical limits of dense retrieval may be reached.
This may cause retrieval to be meaningless and provide no help in answering a query.

\paragraph{MLLM Perception}
This project is based on the assumption that retrieved images provided to a MLLM will further ground generation to visual evidence,
hence conditioning generation on an extra dimension of ground truth information will allow the rate of hallucination and factual errors to be reduced.
Current MLLMs demonstrate expert-level performance on domain-specific knowledge-intensive questions;
but also show deficits in basic visual perception without textual descriptions as a fallback \cite{chen2026babyvisionvisualreasoninglanguage}.
The paper finds that leading MLLMs lack behind adult humans on beyond-language visual tasks by up to 44.4\% (best-performing MLLM 49.7\%; adult humans 94.1\%).
This reflects a severe lack in MLLM visual perception capabilities when the item concerned cannot easily be translated or flattened into natural language.

Our pipeline assumes that fine-grained visual features will assist in answering vision-specific questions that cannot be covered by retrieved text.
However, this assumption is largely limited by MLLM visual perception capabilities.
\vspace{2mm}

\noindent
These limitations motivate the future directions described below.


\subsection{Future Work}

\paragraph{Better Preprocessing Efficiency}
Future work may explore more efficient methods for document preprocessing and inline image captioning to strike a better balance 
between the time taken and captioning accuracy.
Current implementation in our pipeline incurs a linear time complexity in the number of documents and images,
as processing is performed sequentially to avoid exceeding rate limits in API call requests.
Parallel processing and calls, for instance, may be an improvement over the current method.

\paragraph{Foundational Models}
We utilize DINOv2, a foundational vision model, and Talk2DINO, a model that maps text representations into the DINOv2 visual representation space.
These models achieve state-of-the-art performance on various tasks and can be directly used in our pipeline.
For vision-critical applications, the foundational models can be further fine-tuned using domain-specific datasets 
to further raise the capabilties of these models, then plugged into our pipeline to be used in multimodal RAG with a MLLM.

\paragraph{Retrieval}
Further retrieval methods can be employed in further work to further increase retrieval effectiveness.
For example, emerging methods with agentic RAG \cite{gao2026scalingcontextsurveymultimodal} allow dynamic retrieval during generation to access additional visual evidence,
as opposed to the current standard one-off pre-generation retrieval in our pipeline, which may be limited in evidence retrieval for question answering.
Rerankers to re-sort the top-$k$ retrieved items may ensure that the most relevant information is provided to the MLLM for grounding generation.
Cross-modal rerankers could reorder the top-$k$ retrieved images from DINOv2 embeddings based on fine-grained visual similarity to the query,
improving precision at the cost of slightly added latency.

\paragraph{Perception Improvement}
Due to the limits of visual perceptive abilities of existing MLLM frameworks, researchers have proposed that a 
new representation space can be used for tokenizing and embedding text or image patches \cite{chen2026babyvisionvisualreasoninglanguage}.
This may allow visual nuances to be better captured by newly developed models, such that generation can attend more to fine-grained visual elements,
further increasing generation accuracy and reducing hallucination or factual errors.

\paragraph{Application Scenarios}
We apply our multimodal RAG on text-to-text and vision-to-text MLLMs only.
In text-to-vision-enabled models, images and text can be used as queries for generating novel images \cite{NEURIPS2023_43a69d14}.
Such tasks may make use of visual details and nuances in retrieved images to a much larger scale than text generation tasks.
Future work may extend our pipeline and visual element-focused embedding and retrieval method to use with image generation models
for better quality and more realistic generated images.